\textbf{Phần 1. Trắc nghiệm lựa chọn}
\setcounterpageref{ex}{0}

\begin{ex} %[3D2V57]
Tính giá trị của phép chia: $6348 : 3$.
\choice
{$2114$}
{$2126$}
{\True $2116$}
{$2106$}
\loigiai{
$6348 : 3 = 2116$.
\textbf{Đáp án của câu hỏi là C}
}
\end{ex}

\begin{ex} %[3D2V57]
Thực hiện phép chia có chữ số 0 ở thương: $8416 : 4$. Kết quả là
\choice
{\True $2104$}
{$214$}
{$2140$}
{$2014$}
\loigiai{
Lần lượt: $8 : 4 = 2$; $4 : 4 = 1$; $1 : 4 = 0$ dư 1; $16 : 4 = 4$. Kết quả là $2104$.
\textbf{Đáp án của câu hỏi là A}
}
\end{ex}

\begin{ex} %[3D2V57]
Một cuộn cáp quang dài $7525\text{ m}$ được chia đều thành 5 cuộn nhỏ. Mỗi cuộn nhỏ dài bao nhiêu mét?
\choice
{$1550\text{ m}$}
{$155\text{ m}$}
{$1515\text{ m}$}
{\True $1505\text{ m}$}
\loigiai{
$7525 : 5 = 1505\text{ (m)}$.
\textbf{Đáp án của câu hỏi là D}
}
\end{ex}

\begin{ex} %[3D2V57]
Tìm số tự nhiên $y$, biết: $y \times 6 = 7848$.
\choice
{$1380$}
{\True $1308$}
{$1306$}
{$1408$}
\loigiai{
$y = 7848 : 6 = 1308$.
\textbf{Đáp án của câu hỏi là B}
}
\end{ex}

\begin{ex} %[3D2V57]
Thực hiện phép chia có dư: $5687 : 7$. Thương và số dư lần lượt là
\choice
{$812$ và dư $5$}
{$811$ và dư $3$}
{\True $812$ và dư $3$}
{$82$ và dư $3$}
\loigiai{
$5687 = 7 \times 812 + 3$ với số dư $3 < 7$.
\textbf{Đáp án của câu hỏi là C}
}
\end{ex}

\begin{ex} %[3D2V57]
Một tổ sản xuất nhận được $4800\text{ kg}$ nguyên liệu chia đều cho 4 nhóm. Mỗi nhóm đã sử dụng $\dfrac{1}{3}$ số nguyên liệu của mình. Mỗi nhóm còn lại bao nhiêu ki-lô-gam nguyên liệu?
\choice
{\True $800\text{ kg}$}
{$400\text{ kg}$}
{$1200\text{ kg}$}
{$600\text{ kg}$}
\loigiai{
Mỗi nhóm nhận: $4800 : 4 = 1200\text{ kg}$. Đã dùng: $1200 : 3 = 400\text{ kg}$. Còn lại: $1200 - 400 = 800\text{ (kg)}$.
\textbf{Đáp án của câu hỏi là A}
}
\end{ex}

\begin{ex} %[3D2V57]
Trong một phép chia có số chia là 8, thương là 1125 và số dư là số dư lớn nhất có thể có. Số bị chia là
\choice
{$9000$}
{$9008$}
{$9006$}
{\True $9007$}
\loigiai{
Số dư lớn nhất là 7. Số bị chia: $8 \times 1125 + 7 = 9000 + 7 = 9007$.
\textbf{Đáp án của câu hỏi là D}
}
\end{ex}

\begin{ex} %[3D2V57]
Một mảnh đất hình chữ nhật có chu vi $6400\text{ m}$. Biết chiều dài gấp 3 lần chiều rộng. Chiều rộng mảnh đất đó là
\choice
{$1600\text{ m}$}
{\True $800\text{ m}$}
{$600\text{ m}$}
{$1000\text{ m}$}
\loigiai{
Nửa chu vi: $6400 : 2 = 3200\text{ m}$. Tổng số phần: $3 + 1 = 4$ phần. Chiều rộng: $3200 : 4 = 800\text{ (m)}$.
\textbf{Đáp án của câu hỏi là B}
}
\end{ex}

\begin{ex} %[3D2V57]
Tính giá trị của biểu thức: $9135 : 9 - 45 \times 2$.
\choice
{$1015$}
{$915$}
{\True $925$}
{$935$}
\loigiai{
$9135 : 9 = 1015$; $45 \times 2 = 90$; $1015 - 90 = 925$.
\textbf{Đáp án của câu hỏi là C}
}
\end{ex}

\begin{ex} %[3D2V57]
Một người chia đều $3690\text{ kg}$ ngô vào 9 bao to. Người đó đã bán đi 4 bao to. Khối lượng ngô còn lại là
\choice
{\True $2050\text{ kg}$}
{$1640\text{ kg}$}
{$2100\text{ kg}$}
{$1950\text{ kg}$}
\loigiai{
Mỗi bao chứa: $3690 : 9 = 410\text{ kg}$. Còn lại 5 bao: $410 \times 5 = 2050\text{ (kg)}$.
\textbf{Đáp án của câu hỏi là A}
}
\end{ex}

\begin{ex} %[3D2V57]
Tìm một số biết rằng nếu lấy số đó chia cho 7 rồi cộng với 125 thì được 1125. Số đó là
\choice
{$6500$}
{$7200$}
{$6800$}
{\True $7000$}
\loigiai{
Số đó chia cho 7 bằng: $1125 - 125 = 1000$. Số cần tìm: $1000 \times 7 = 7000$.
\textbf{Đáp án của câu hỏi là D}
}
\end{ex}

\begin{ex} %[3D2V57]
Người ta xếp 4520 viên gạch lên 4 xe bò đều nhau. Mỗi xe chở được bao nhiêu viên gạch?
\choice
{$1120$ viên}
{\True $1130$ viên}
{$1140$ viên}
{$1150$ viên}
\loigiai{
$4520 : 4 = 1130$ viên gạch.
\textbf{Đáp án của câu hỏi là B}
}
\end{ex}

\textbf{Phần 2. Trắc nghiệm đúng sai}
\setcounterpageref{ex}{0}

\begin{ex} %[3D2V57]
Xét tính đúng sai của phép chia số có bốn chữ số cho số có một chữ số:
\choiceTF
{\True Phép chia thực hiện từ trái sang phải (từ hàng cao nhất).}
{\True $6024 : 6 = 1004$.}
{$6024 : 6 = 14$.}
{Trong phép chia cho 5, số dư có thể bằng 5.}
\loigiai{
\begin{itemize}
    \item Ý a) Đúng: Quy tắc chia chuẩn.
    \item Ý b) Đúng: Chia có hai chữ số 0 ở thương.
    \item Ý c) Sai: Sai vì bỏ qua các chữ số 0 ở thương.
    \item Ý d) Sai: Sai vì số dư luôn nhỏ hơn số chia.
\end{itemize}
\begin{itemize}
    \item a) \textbf{Đúng}.
    \item b) \textbf{Đúng}.
    \item c) \textbf{Sai}.
    \item d) \textbf{Sai}.
\end{itemize}
}
\end{ex}

\begin{ex} %[3D2V57]
Một nhà máy có 5430 sản phẩm cần đóng vào các thùng loại 6 sản phẩm:
\choiceTF
{\True Đóng được tất cả 905 thùng.}
{\True Không còn thừa sản phẩm nào sau khi đóng.}
{Nếu mỗi xe chở được 100 thùng thì cần 8 xe để chở hết số thùng trên.}
{\True Năm thùng chứa tất cả $6 \times 5 = 30$ sản phẩm.}
\loigiai{
\begin{itemize}
    \item Ý a) Đúng: $5430 : 6 = 905$ thùng.
    \item Ý b) Đúng: 5430 chia hết cho 6.
    \item Ý c) Sai: Sai vì cần ít nhất 10 xe ($905 : 100 = 9$ dư 5).
    \item Ý d) Đúng: Đúng với quy cách đóng gói.
\end{itemize}
\begin{itemize}
    \item a) \textbf{Đúng}.
    \item b) \textbf{Đúng}.
    \item c) \textbf{Sai}.
    \item d) \textbf{Đúng}.
\end{itemize}
}
\end{ex}

\textbf{Phần 3. Trả lời ngắn}
\setcounterpageref{ex}{0}

\begin{ex} %[3D2V57]
Tính kết quả phép chia: $8472 : 8$.
\shortans{1059}
\loigiai{
\textbf{Đáp số:} $1059$\\
$8472 : 8 = 1059$.
}
\end{ex}

\begin{ex} %[3D2V57]
Tìm số tự nhiên $x$, biết: $x \times 7 = 8421$.
\shortans{1203}
\loigiai{
\textbf{Đáp số:} $1203$\\
$x = 8421 : 7 = 1203$.
}
\end{ex}

\begin{ex} %[3D2V57]
Một mảnh đất hình vuông có chu vi $6840\text{ m}$. Tính độ dài cạnh của mảnh đất theo mét.
\shortans{1710}
\loigiai{
\textbf{Đáp số:} $1710$\\
$6840 : 4 = 1710\text{ (m)}$.
}
\end{ex}

\begin{ex} %[3D2V57]
Người ta chia đều $7245\text{ kg}$ gạo cho 5 cửa hàng. Hỏi mỗi cửa hàng nhận được bao nhiêu ki-lô-gam gạo?
\shortans{1449}
\loigiai{
\textbf{Đáp số:} $1449$\\
$7245 : 5 = 1449\text{ (kg)}$.
}
\end{ex}

\begin{ex} %[3D2V57]
Tìm số dư trong phép chia: $4589 : 9$.
\shortans{8}
\loigiai{
\textbf{Đáp số:} $8$\\
$4589 = 9 \times 509 + 8$. Số dư là 8.
}
\end{ex}

\begin{ex} %[3D2V57]
Một trang trại có 3696 con gà chia đều vào 6 chuồng. Sau đó người ta bắt đi 120 con ở chuồng thứ nhất. Hỏi chuồng thứ nhất còn lại bao nhiêu con gà?
\shortans{496}
\loigiai{
\textbf{Đáp số:} $496$\\
Mỗi chuồng ban đầu: $3696 : 6 = 616$ con. Còn lại: $616 - 120 = 496$ con.
}
\end{ex}
